\documentclass[11pt]{article}
\usepackage[margin=0.9in]{geometry}
\usepackage{amsmath,amssymb,amsthm}
\usepackage[T1]{fontenc}
\usepackage{lmodern,microtype}
\usepackage[colorlinks=true,linkcolor=blue,urlcolor=blue]{hyperref}
\newtheorem{theorem}{Theorem}[section]
\newtheorem{lemma}[theorem]{Lemma}
\newtheorem{proposition}[theorem]{Proposition}
\newtheorem{corollary}[theorem]{Corollary}
\theoremstyle{remark}
\newtheorem{remark}[theorem]{Remark}
\newcommand{\norm}[1]{\left\|#1\right\|}
\newcommand{\diag}{\operatorname{diag}}
\newcommand{\ran}{\operatorname{ran}}
\newcommand{\R}{\mathbb R}
\newcommand{\C}{\mathbb C}
\title{A uniform flat auxiliary-current neighborhood}
\author{Notebook 09 integration fragment}
\date{16 October 2026}
\begin{document}
\maketitle

\noindent\textbf{Integration note.}
The body below follows the full flat inverse theorem and the scalar
source-gauge locality lemma. Those antecedents have separate internal
derivations and audits. This fragment supplies the current argument;
it does not replace their proofs. Its source application uses the
deposited homogeneous flat operator definitions.

% BEGIN NOTEBOOK FRAGMENT
\section{A uniform neighborhood for the flat auxiliary current}
\label{current-section}

We keep the background equal to the identity and vary the current
independently. Fine spacing is $M^{-1}$, with odd $M$, and the fine
measure is $M^{-d}$. The one-form average is the corrected source
average at flatness, already identified with $Q=Q_M^{\mathrm{st}}$.
The scalar source projection is $R$, and $D$ is the fine scalar
gradient. Write
\[
 \Delta_0=d_\eta^*d_\eta+DRD^*+Q^*aQ,
 \qquad G_0=\Delta_0^{-1},
 \qquad q_0=(QG_0Q^*)^{-1}.
\]
The homogeneous source convention has $a=1$ at unit physical block
scale. All conclusions below also hold uniformly for a fixed compact
interval $0<a_-\le a\le a_+<\infty$. We do not assert a bound for an
unmatched sequence of penalties.

The preceding flat inverse theorem and minimizing-extension theorem give
\begin{equation}
 q_0=q_M+aI,
 \qquad E_0=G_0Q^*q_0,
 \qquad QE_0=I,
 \qquad RD^*E_0=0.
 \label{current-base-identities}
\end{equation}
Here $q_M$ is the unpenalized constrained Maxwell precision. The
addition of $aI$ must be retained in $q_0$. The normalized extension
$E_0$ is unchanged by this penalty; a source convention
$QH_0=\ell^{-1}$ instead uses $H_0=E_0\ell^{-1}$.

All adjoints initially use these physical measures at the real flat
reference. At complex current they denote the fixed bilinear
transpose. They do not conjugate the varying complex current.

\subsection{Current vertices and the weighted algebra}
\label{current-vertices-section}

For a fine bond $b$, the flat gauge expansion has mixed and polarized
quadratic terms
\begin{align}
 c^{(1)}(A,\lambda)(b)
 &=\frac{i}{2}[\lambda(b_-)+\lambda(b_+),A(b)],
 \label{current-c1}\\
 2c^{(2)}(\lambda_1,\lambda_2)(b)
 &=\frac{1}{2i}\bigl(
 [\lambda_1(b_-),(D\lambda_2)(b)]
 +[\lambda_2(b_-),(D\lambda_1)(b)]\bigr).
 \label{current-c2}
\end{align}
These are the deposited formulas at Raw I:13698--13718 and
I:13811--13821. Define the current-linear maps by
\[
 \langle\lambda,\mathfrak j_J A\rangle
   =\langle c^{(1)}(A,\lambda),J\rangle,
 \qquad \mathfrak k_J=\mathfrak j_J^*,
\]
\[
 \langle\lambda_1,\mathfrak m_J\lambda_2\rangle
   =2\langle c^{(2)}(\lambda_1,\lambda_2),J\rangle.
\]
All bond and site pairings carry the same fine volume weight.
Set $\Lambda=G'RD^*$, with the fixed scalar source inverse $G'$.
The scalar locality lemma gives uniform weighted bounds for
$\Lambda$ and $D\Lambda$. The auxiliary perturbation is
\begin{equation}
 \Theta'_J=\mathfrak k_J\Lambda+
       \Lambda^*\mathfrak j_J+\Lambda^*\mathfrak m_J\Lambda.
 \label{current-theta-definition}
\end{equation}
This is precisely the auxiliary perturbation in Raw I:13725--13731.
It is distinct from the additional nonlinear-constraint operator
denoted there by $\Theta_J$.

Group fine fields into unit coarse cells with their physical Hilbert
norms. For block operators between these fine fibers, or between
fine and fixed-dimensional coarse fibers, use
\begin{equation}
 \norm{T}_\kappa=
 \max\left\{
   \sup_x\sum_y e^{\kappa d(x,y)}\norm{T_{xy}},
   \sup_y\sum_x e^{\kappa d(x,y)}\norm{T_{xy}}
 \right\}.
 \label{current-schur-norm}
\end{equation}
Here $d(x,y)=|x-y|_1$ in infinite volume and is periodic coarse-cell
distance on a finite compatible torus. This norm is submultiplicative,
including rectangular products, and transpose preserves it.

\begin{lemma}[A current bound without derivatives of the current]
\label{current-theta-bound}
There is $\theta<\infty$, independent of refinement and compatible
periodic volume, such that
\begin{equation}
 \norm{\Theta'_J}_\kappa\le\theta\norm{J}_\infty.
 \label{current-theta-norm}
\end{equation}
The map $J\mapsto\Theta'_J$ is complex linear. The current requires no
regularity beyond boundedness.
\end{lemma}

\begin{proof}
Define the finite-incidence map $V_J$ from fine bonds to fine sites by
\[
 \langle\lambda,V_J A\rangle
   =\frac{1}{2i}\langle[\lambda(b_-),A(b)],J(b)\rangle.
\]
Equation \eqref{current-c2} gives the exact form identity
\[
 \mathfrak m_J=V_JD+D^*V_J^*.
\]
Thus \eqref{current-theta-definition} becomes
\begin{equation}
 \Theta'_J=
 \mathfrak j_J^*\Lambda+\Lambda^*\mathfrak j_J
 +\Lambda^*V_JD\Lambda+(D\Lambda)^*V_J^*\Lambda.
 \label{current-theta-factorization}
\end{equation}
If $|[X,Y]|\le c_{\rm br}|X||Y|$, both
$\mathfrak j_J$ and $V_J$ have fine Hilbert operator norm at most
$C_d c_{\rm br}\norm{J}_\infty$. Indeed there are at most $2d$
incident bonds per site and two endpoints per bond; Cauchy--Schwarz
and the common physical volume weight give this bound without a
power of $M$. These coefficient maps have a fixed finite cell range,
so the same estimate holds in \eqref{current-schur-norm}.

Apply the weighted algebra inequality to
\eqref{current-theta-factorization}, using the established bounds for
$\Lambda,D\Lambda$. Every derivative is on a scalar leg through
$D\Lambda$. No derivative acts on $J$, even if $J$ jumps at a cell
boundary. The proof also applies to complex currents with the fixed
bilinear transpose.
\end{proof}

Choose a common sufficiently small $\kappa>0$. The flat inverse,
averaging and Maxwell locality results give constants $g,b,q$ with
\begin{equation}
 \norm{G_0}_\kappa\le g,
 \qquad \norm{Q}_\kappa,\norm{Q^*}_\kappa\le b,
 \qquad \norm{q_0}_\kappa\le q.
 \label{current-input-bounds}
\end{equation}
The fine inverse result also gives all physical fine-difference
sandwiches of total order at most two. These are established
homogeneous flat inputs; they are not hypotheses about a nonflat
source carrier. None of these constants includes a factor for the
number of sites in a fine cell.

\subsection{Uniform inverses and the exact quadratic remainder}
\label{current-neighborhood-section}

Define
\begin{equation}
 c=gb^2q,\qquad e=gbq,\qquad t_0=g(1+c),\qquad
 \varepsilon=\frac{1}{4g\theta(1+c)}.
 \label{current-radius}
\end{equation}
If $\theta=0$ the perturbation vanishes and the current restriction is
unnecessary. Otherwise use the displayed positive radius.

\begin{theorem}[Uniform flat auxiliary-current neighborhood]
\label{current-neighborhood-theorem}
For every complex bounded current with
$\norm{J}_\infty<\varepsilon$, the operators
\begin{equation}
 G_J=(\Delta_0-\Theta'_J)^{-1},\qquad
 q_J=(QG_JQ^*)^{-1},\qquad E_J=G_JQ^*q_J
 \label{current-family}
\end{equation}
exist and depend holomorphically on $J$ in the weighted algebra. They
are the deposited auxiliary family, with
$G_J=G_\pi^J$, $q_J=\mathsf q_J$ and $E_J=H^J\ell$. Uniformly in odd $M$, compatible volume, and
$a\in[a_-,a_+]$,
\begin{equation}
 \norm{G_J}_\kappa\le2g,\qquad
 \norm{q_J}_\kappa\le2q,\qquad
 QE_J=I,\qquad\norm{E_J}_\kappa\le4e.
 \label{current-uniform-bounds}
\end{equation}
Let
\begin{equation}
 B_0=G_0-E_0QG_0=G_0-E_0q_0^{-1}E_0^*.
 \label{current-B0}
\end{equation}
Then the exact formulas are
\begin{align}
 E_J-E_0
   &=(I-B_0\Theta'_J)^{-1}B_0\Theta'_JE_0,
 \label{current-extension-remainder}\\
 q_J-q_0
   &=-E_0^*\Theta'_J(I-B_0\Theta'_J)^{-1}
                  B_0\Theta'_JE_0.
 \label{current-quadratic-remainder}
\end{align}
In particular,
\begin{equation}
 \norm{E_J-E_0}_\kappa\le2t_0\theta e\norm{J}_\infty,
 \qquad
 \norm{q_J-q_0}_\kappa\le
       2e^2\theta^2t_0\norm{J}_\infty^2.
 \label{current-remainder-bounds}
\end{equation}
\end{theorem}

\begin{proof}
Abbreviate $\Theta=\Theta'_J$. The first Neumann series gives
\[
 G_J=(I-G_0\Theta)^{-1}G_0,
 \qquad
 \norm{G_J-G_0}_\kappa\le2g^2\theta\norm{J}_\infty,
\]
because $\norm{G_0\Theta}_\kappa<1/4$.
Writing $W_J=QG_JQ^*$, we obtain
\[
 \norm{q_0(W_J-W_0)}_\kappa
 \le2qb^2g^2\theta\norm{J}_\infty
 <\frac{c}{2(1+c)}<\frac12.
\]
The second Neumann series, with its specified order, is
\[
 q_J=[I+q_0(W_J-W_0)]^{-1}q_0.
\]
This proves \eqref{current-uniform-bounds}. Both series converge in
the same Banach algebra and therefore prove holomorphy on the entire
stated complex ball.

The constrained covariance has
\[
 QB_0=B_0Q^*=0,\qquad
 B_0\Delta_0=I-E_0Q,\qquad
 \norm{B_0}_\kappa\le t_0.
\]
Apply $B_0$ to
$(\Delta_0-\Theta)E_J=Q^*q_J$ and use $QE_J=I$. This gives
$E_J-E_0=B_0\Theta E_J$, hence
\eqref{current-extension-remainder}. No commutation of $B_0$ with
$\Theta$ is used. The chosen radius has
$\norm{B_0\Theta}_\kappa<1/4$.

Multiplying the stationary equation by $E_0^*$ gives
$q_J=q_0-E_0^*\Theta E_J$. The flat source gauge implies
$\Lambda E_0=0$, and \eqref{current-theta-definition} then gives
\begin{equation}
 \Theta'_JE_0=\Lambda^*\mathfrak j_JE_0,
 \qquad E_0^*\Theta'_JE_0=0.
 \label{current-linear-cancellation}
\end{equation}
Substitution of \eqref{current-extension-remainder} proves
\eqref{current-quadratic-remainder}. Finally
$\norm{E_0}_\kappa\le e$ and
$\norm{(I-B_0\Theta)^{-1}}_\kappa\le2$ give
\eqref{current-remainder-bounds}.
\end{proof}

\subsection{Factorial derivatives and localized current legs}
\label{current-derivatives-section}

For a current in this ball put
\begin{equation}
 B_J=G_J-E_JQG_J,
 \qquad t=2g(1+4c).
 \label{current-BJ}
\end{equation}
Then $B_J$ is the symmetric constrained covariance and
$\norm{B_J}_\kappa\le t$. All directional derivatives below are
derivatives in current space, not fine spatial differences.

\begin{proposition}[Exact derivative chains]
\label{current-derivative-chains}
For $n\ge1$ and bounded current directions $J_1,\ldots,J_n$,
\begin{equation}
 \begin{split}
 D^nq_J[J_1,\ldots,J_n]
 =-\sum_{\sigma\in S_n}
 E_J^*\Theta'_{J_{\sigma(1)}}B_J
 \Theta'_{J_{\sigma(2)}}\cdots B_J
 \Theta'_{J_{\sigma(n)}}E_J.
 \end{split}
 \label{current-q-derivative-chain}
\end{equation}
Each summand contains $n-1$ factors $B_J$; for $n=1$ it is
$-E_J^*\Theta'_{J_1}E_J$. The corresponding formula for $D^nE_J$
is the sum of the ordered products
$B_J\Theta'_{J_{\sigma(1)}}\cdots
B_J\Theta'_{J_{\sigma(n)}}E_J$. Consequently
\begin{align}
 \norm{D^nq_J[J_1,\ldots,J_n]}_\kappa
 &\le n!(4e)^2\theta^n t^{n-1}
                   \prod_i\norm{J_i}_\infty,
 \label{current-q-factorial}\\
 \norm{D^nE_J[J_1,\ldots,J_n]}_\kappa
 &\le n!(4e)(t\theta)^n\prod_i\norm{J_i}_\infty.
 \label{current-E-factorial}
\end{align}
\end{proposition}

\begin{proof}
At a fixed $J$, repeat the stationary-equation argument for a
sufficiently small increment $K$. It gives the exact identities
\begin{align}
 B_{J+K}&=(I-B_J\Theta'_K)^{-1}B_J,
 \label{current-B-increment}\\
 E_{J+K}&=(I-B_J\Theta'_K)^{-1}E_J,
 \label{current-E-increment}\\
 q_{J+K}&=q_J-E_J^*\Theta'_K
                    (I-B_J\Theta'_K)^{-1}E_J.
 \label{current-q-increment}
\end{align}
The first is the resolvent identity for the quadratic form
restricted to $\ker Q$, with its inverse extended by zero on
$\ran Q^*$. Expanding the series and taking the coefficient
multilinear in the directions produces exactly the permutations in
\eqref{current-q-derivative-chain} and the extension formula. The
norm bounds follow from submultiplicativity. The argument is valid
at each point in the ball by using a sufficiently small increment;
it does not enlarge the current domain.
\end{proof}

At zero current $Dq_0[J_1]=0$ by
\eqref{current-linear-cancellation}. The earlier scalar gauge
identity, with $Z=G'RG'$, gives the more explicit specializations
\[
 DE_0[J_1]=DZ\mathfrak j_{J_1}E_0,
\]
\begin{equation}
 \begin{split}
 D^2q_0[J_1,J_2]
 ={}&-(\mathfrak j_{J_1}E_0)^*Z(\mathfrak j_{J_2}E_0)\\
    &-(\mathfrak j_{J_2}E_0)^*Z(\mathfrak j_{J_1}E_0).
 \end{split}
 \label{current-zero-Hessian}
\end{equation}
The first extension response is a fine gradient. It need not
preserve the original source gauge at nonzero current.

\begin{proposition}[Paths through the current supports]
\label{current-localized-paths}
Let $J_{x_i}$ be supported on fine bonds based in cell $x_i$.
There are constants $C_0,C_1$, independent of refinement, compatible
volume, source locations and $n$, such that
\begin{equation}
 \begin{split}
 &\norm{D^nq_J[J_{x_1},\ldots,J_{x_n}](b,c)}\\
 &\quad\le C_0 C_1^n\prod_i\norm{J_{x_i}}_\infty
 \sum_{\sigma\in S_n}
 \exp\!\left[-\kappa\left(
 d(b,x_{\sigma(1)})
 +\sum_{i=1}^{n-1}d(x_{\sigma(i)},x_{\sigma(i+1)})
 +d(x_{\sigma(n)},c)\right)\right].
 \end{split}
 \label{current-path-bound}
\end{equation}
Here $b,c$ are coarse cells with their fixed orientation/color
fibers. Repeated source cells are allowed. The extension
derivatives have the same path structure, with the appropriate
fine/coarse end fibers.
\end{proposition}

\begin{proof}
For a current supported in cell $x$, each coefficient vertex
$\mathfrak j_{J_x}$ or $V_{J_x}$ reads and writes only a fixed
neighboring-cell set of $x$, and its restricted operator norm is
$O(\norm{J_x}_\infty)$. Expand each perturbation in a chain using
\eqref{current-theta-factorization}. Between consecutive local
vertices are products of the form $A B_J C^*$, with the appropriate
scalar/bond domains and with $A,C$ among
$I,\Lambda,D\Lambda$. Such products have a uniform weighted Schur
bound and hence their $(x,y)$ block is bounded by a constant times
$e^{-\kappa d(x,y)}$. The end products use $A E_J$ and obey the
same statement.

There are at most four choices per perturbation and a fixed number
of neighboring cells per coefficient vertex. Their sizes and
bounded offsets are absorbed into $C_1^n$. Multiplying the block
bounds and then summing the permutations proves
\eqref{current-path-bound}. This uses physical block operator norms
before multiplying; it never counts the fine sites in a block.
It also avoids convolution of pointwise bounds at a borderline
exponential rate. Periodic distance obeys the same triangle
inequality, so the proof is unchanged on compatible tori.
\end{proof}

\subsection{Fine derivatives, finite volume, and the scope of the result}
\label{current-scope-section}

The proved full flat inverse theorem controls every specified fine
difference sandwich of total order at most two. Write
$\mathcal D_\alpha$ and $\mathcal D_\beta$ for fixed words of those
orders, with the appropriate orientation domains. The exact identity
\begin{equation}
 \begin{split}
 \mathcal D_\alpha G_J\mathcal D_\beta
 ={}&\mathcal D_\alpha G_0\mathcal D_\beta\\
 &+(\mathcal D_\alpha G_0)\Theta'_J
       (I-G_0\Theta'_J)^{-1}G_0\mathcal D_\beta
 \end{split}
 \label{current-spatial-derivative-transfer}
\end{equation}
therefore gives the same uniform weighted bounds when
$\alpha+\beta\le2$. Derivatives act on the outside propagator
factors, not on the current. The one-sided estimates also give
the corresponding output fine derivatives of $E_J$. These are
physical cell $L^2$ operator estimates; no pointwise or H\"older
carrier estimate is inferred from them.

The reference kernels on finite compatible tori are exact
periodizations of the homogeneous kernels. All subsequent arguments
take place in the weighted block algebra on that torus. They apply
to arbitrary bounded periodic currents, without requiring them to
be constant from cell to cell, and retain constants uniform even
for short periods. The infinite-volume formulas follow in the
same algebra. No passage through an unweighted growing-dimensional
matrix norm is used.

The result is a uniform complex neighborhood for the displayed
\emph{auxiliary} flat family, with a quadratic precision remainder,
factorial current derivatives, and paths through their current
supports. The full source inverse instead has
\[
 G_1^J=[\Delta_0-\Theta'_J-\Pi^*\Theta_J\Pi]^{-1},
 \qquad \Pi=I-D\Lambda.
\]
The separate nonlinear-constraint term $\Theta_J$ is not covered
here. Its first coarse-precision derivative can already contain
$-E_0^*\Theta^{[1]}_J E_0$. Thus the cancellation
\eqref{current-linear-cancellation} is not a statement that the
full physical precision has zero current derivative. Nonflat
backgrounds, retained multiscale boundaries, source pointwise
carrier norms, and the interacting ultraviolet limit remain
outside this theorem.
% END NOTEBOOK FRAGMENT
\end{document}
